\section{Background and Terminology}
\label{sec:background}

% 本节的任务是【固定术语】，不是做背景科普。
% 全文必须对 harness / scaffold / harness self-evolution 三个词有唯一定义。

\subsection{Recursive self-improvement}

\begin{quote}
A system exhibits recursive self-improvement if the improved system can
participate in producing the next improvement, forming a cycle of
evaluate--update--re-evaluate.
\end{quote}

This definition, adopted from the course framing \todo{confirm exact
wording}, has a consequence that is easy to overlook. \emph{Producing a better
artifact once} is not recursive improvement. What makes it recursive is the
participation of the improved artifact in the next round. We return to this in
Section~\ref{sec:taxonomy}, where it becomes the final axis of our taxonomy.

\subsection{Three loci of self-improvement}

Self-improving AI systems are commonly described along three loci:

\begin{table}[H]
\centering
\caption{The three loci of self-improvement and their targets.}
\label{tab:loci}
\begin{tabularx}{\textwidth}{@{}l X l@{}}
\toprule
\thead{Locus} & \thead{What changes} & \thead{In this review} \\
\midrule
Model      & Parameters and the training procedure & No \\
Data       & Datasets, experience, and their production pipeline & No \\
\textbf{Harness} & \textbf{The program that organizes model, data, and environment interaction} & \textbf{Yes} \\
\bottomrule
\end{tabularx}
\end{table}

The harness is the outermost of the three: it does not touch parameters or
data, but determines how both are put to use. \todo{Add one or two sentences on
why this boundary matters for the present review.}

\subsection{What counts as a harness}

We adopt the following working definition:

\begin{quote}
A \emph{harness} is the executable program that mediates an agent's
interaction with its environment --- comprising, at minimum, its prompt and
context construction, its tool interface, its control flow and termination
logic, and any persistent memory or evaluator it consults.
\end{quote}

\begin{quote}
A system \emph{self-evolves its harness} when the program is the target of its
own modification, rather than merely the medium through which a fixed
improvement procedure acts.
\end{quote}

This definition draws on the scaffold framing of \cite{zelikman2023stop} and
the operational-scaffold framing of \cite{ren2026selfimprovements}, which
decomposes an agent into a foundation model coupled with prompts, memory,
tools, and control logic.

\subsection{Terminological cautions}

\begin{itemize}
    \item \textbf{Harness} and \textbf{scaffold} are used interchangeably in
    this literature. We use \emph{harness} throughout for consistency.

    \item \textbf{Recursion} is claimed loosely. Authors differ on whether
    modifying only the artifact, or also the procedure that produced it,
    qualifies \todo{quote STOP's own words here as the worked example; see
    Section~\ref{sec:open}}.

    \item \textbf{Self-improvement} is claimed with varying evidentiary
    standards. We therefore record, for each paper, whether a compute-matched
    baseline was used and whether search cost was reported.
\end{itemize}

% ----------------------------------------------------------------------
% 候选引用（核对后启用）:
%   \cite{zelikman2023stop}          STOP --- harness 概念源头
%   \cite{hu2024adas}                ADAS --- 架构搜索
%   \cite{robeyns2025sica}           SICA --- 自编辑代码
%   \cite{zhang2025dgm}              DGM --- 完整 RSI 声明
%   \cite{novikov2025alphaevolve}    AlphaEvolve --- 工业界
%   \cite{ning2026codeharness}       Code as Agent Harness --- 概念框架
%   \cite{ren2026selfimprovements}   Survey --- operational scaffold
%   \cite{fang2025comprehensive}     Survey --- 自进化 agent 全景
% ----------------------------------------------------------------------
